
\documentclass[11pt,a4paper]{article}

% ------------------------------------------------
% Packages
% ------------------------------------------------
\usepackage[utf8]{inputenc}
\usepackage[T1]{fontenc}
\usepackage{lmodern}

\usepackage{amsmath,amssymb,amsthm}
\usepackage{mathtools}
\usepackage{geometry}
\usepackage{hyperref}
\usepackage{microtype}
\usepackage{enumitem}

% ------------------------------------------------
% Paragraph Formatting
% ------------------------------------------------
\setlength{\parindent}{0pt}
\setlength{\parskip}{6pt}

\geometry{margin=1in}

% ------------------------------------------------
% Theorem Environments
% ------------------------------------------------
\newtheorem{theorem}{Theorem}[section]
\newtheorem{proposition}[theorem]{Proposition}
\newtheorem{lemma}[theorem]{Lemma}

\theoremstyle{definition}
\newtheorem{example}[theorem]{Example}
\newtheorem{definition}[theorem]{Definition}

\theoremstyle{remark}
\newtheorem{remark}[theorem]{Remark}

% ------------------------------------------------
% Title
% ------------------------------------------------
\title{\textbf{An Introduction to Feynman's
Technique for Parametric Integration}\\[0.3cm]
\large Applications and Mathematical Justification}

\author{Biajid E Chowdhury}
\date{\today}

\begin{document}

\maketitle


% =================================================
\section{Introduction}
% =================================================

Integration is one of the central operations of
calculus. Although elementary techniques such as
substitution, integration by parts, and partial
fraction decomposition allow us to evaluate many
integrals, there are numerous integrals arising
in physics and applied mathematics for which
these methods are either inefficient or do not
lead to an elementary antiderivative.

In some situations, however, an apparently
difficult integral becomes considerably easier
when we introduce an additional parameter and
study how the resulting integral changes with
respect to that parameter.

This observation is the central idea behind
\emph{parametric integration}, commonly known
as \emph{Feynman's technique of differentiating
under the integral sign}.

The method is particularly useful when
differentiating the integrand with respect to
a parameter simplifies its algebraic structure.

To illustrate the general idea, suppose that
we wish to evaluate an integral of the form

\begin{equation}
    I=\int_a^b f(x)\,dx.
    \label{eq:original_integral}
\end{equation}

Instead of attempting to evaluate
\eqref{eq:original_integral} directly, we
introduce a parameter $\alpha$ and construct
a family of integrals

\begin{equation}
    F(\alpha)
    =
    \int_a^b f(x,\alpha)\,dx,
    \label{eq:parametric_integral}
\end{equation}

where the original integral is recovered at
some particular value $\alpha=\alpha_0$:

\[
    I=F(\alpha_0).
\]

The key observation is that, under suitable
mathematical conditions, differentiation
with respect to $\alpha$ can be performed
inside the integral:

\begin{equation}
    \boxed{
    \frac{d}{d\alpha}
    \int_a^b f(x,\alpha)\,dx
    =
    \int_a^b
    \frac{\partial f(x,\alpha)}
         {\partial\alpha}\,dx
    }
    \label{eq:feynman_rule}
\end{equation}

The expression on the right-hand side may
be substantially easier to evaluate than
the original integral.

Once we determine $F'(\alpha)$, we can
recover $F(\alpha)$ by integration with
respect to $\alpha$. Any resulting
constant of integration can then be
determined from a known value of
$F(\alpha)$ at a convenient parameter,
or through a suitable limiting argument.

Finally, evaluating $F$ at $\alpha_0$
recovers the desired integral.

The essential strategy is therefore not
to attack a difficult integral directly,
but to embed it in a family of integrals
whose dependence on a parameter is
easier to analyze.

\medskip

However, an important mathematical question
remains. The identity \eqref{eq:feynman_rule}
is not universally valid for arbitrary
functions. Differentiation and integration
cannot always be interchanged.

Why, then, does Feynman's technique work?

We will postpone this question until
after developing several applications.
We will then investigate the
mathematical justification through
Leibniz's integral rule and, finally,
interpret the technique using the
language of linear operators.



% =================================================
\section{Applications of Feynman's Technique}
% =================================================

In this section, we will solve several
integral problems using Feynman's technique
of differentiating under the integral sign.
Our examples will include Gaussian integrals,
trigonometric integrals, and logarithmic
integrals, some of which are difficult
to evaluate using conventional methods.

For the moment, we ask the reader to
concentrate on the mechanics of the method:
how to introduce an appropriate parameter,
how differentiation with respect to that
parameter simplifies the integrand, and
how the resulting expression leads to
the value of the original integral.

We will temporarily assume that
differentiation and integration can be
interchanged whenever we apply the method.
The mathematical conditions under which
this interchange is justified will be
discussed in a later section.

Our immediate objective is to develop
an intuitive and practical understanding
of the technique through explicit
calculations.

% -------------------------------------------------
\subsection{Gaussian Integrals: Odd Powers}
% -------------------------------------------------

Let us begin by considering the family
of Gaussian integrals

\begin{equation}
    I_n =
    \int_0^\infty x^n e^{-x^2}\,dx,
    \qquad n=0,1,2,\ldots
    \label{eq:gaussian_general}
\end{equation}

These integrals arise frequently in
probability theory, statistics, and
mathematical physics.

We will treat odd and even values of
$n$ separately, since the corresponding
calculations require slightly different
approaches.

First, suppose that $n$ is odd.
We may write

\[
    n=2m+1,
    \qquad m=0,1,2,\ldots
\]

Our objective is therefore to evaluate

\begin{equation}
    I_{2m+1}
    =
    \int_0^\infty
    x^{2m+1}e^{-x^2}\,dx.
    \label{eq:gaussian_odd}
\end{equation}

\subsubsection*{Step 1: A preliminary substitution}

Before introducing a parameter, we
make the substitution

\[
    u=x^2,
    \qquad du=2x\,dx.
\]

Since

\[
    x^{2m+1}\,dx
    =
    (x^2)^m x\,dx
    =
    \frac{1}{2}u^m\,du,
\]

equation \eqref{eq:gaussian_odd}
becomes

\begin{equation}
    I_{2m+1}
    =
    \frac{1}{2}
    \int_0^\infty u^m e^{-u}\,du.
    \label{eq:gaussian_odd_substitution}
\end{equation}

Although the substitution has
simplified the exponential,
we still have the factor $u^m$
inside the integral.

This is where Feynman's technique
becomes useful.

\subsubsection*{Step 2: Introducing a parameter}

Consider the parameter-dependent
integral

\begin{equation}
    F(\alpha)
    =
    \int_0^\infty e^{-\alpha u}\,du,
    \qquad \alpha>0.
    \label{eq:gaussian_parameter}
\end{equation}

The reason for introducing the
parameter in the exponent is that

\[
    \frac{\partial}{\partial\alpha}
    e^{-\alpha u}
    =
    -u e^{-\alpha u}.
\]

Differentiating a second time gives

\[
    \frac{\partial^2}{\partial\alpha^2}
    e^{-\alpha u}
    =
    u^2 e^{-\alpha u}.
\]

More generally,

\begin{equation}
    \frac{\partial^m}{\partial\alpha^m}
    e^{-\alpha u}
    =
    (-1)^m u^m e^{-\alpha u}.
    \label{eq:gaussian_parameter_derivative}
\end{equation}

Thus, repeated differentiation
with respect to the parameter
generates precisely the powers
of $u$ that appear in our
original integral.

\subsubsection*{Step 3: Applying Feynman's technique}

Differentiating
\eqref{eq:gaussian_parameter}
$m$ times with respect to
$\alpha$, we obtain

\begin{align}
    F^{(m)}(\alpha)
    &=
    \frac{d^m}{d\alpha^m}
    \int_0^\infty e^{-\alpha u}\,du
    \notag\\
    &=
    \int_0^\infty
    \frac{\partial^m}{\partial\alpha^m}
    e^{-\alpha u}\,du
    \notag\\
    &=
    (-1)^m
    \int_0^\infty
    u^m e^{-\alpha u}\,du.
    \label{eq:gaussian_odd_differentiation}
\end{align}

Consequently,

\begin{equation}
    \int_0^\infty
    u^m e^{-\alpha u}\,du
    =
    (-1)^m F^{(m)}(\alpha).
    \label{eq:gaussian_odd_identity}
\end{equation}

Comparing this expression with
\eqref{eq:gaussian_odd_substitution},
we see that the desired integral
can be recovered by setting
$\alpha=1$.

Therefore,

\begin{equation}
    I_{2m+1}
    =
    \frac{(-1)^m}{2}
    \left.
    F^{(m)}(\alpha)
    \right|_{\alpha=1}.
    \label{eq:gaussian_odd_recovery}
\end{equation}

\subsubsection*{Step 4: Evaluating the parametric integral}

The integral defining $F(\alpha)$
is elementary:

\begin{align}
    F(\alpha)
    &=
    \int_0^\infty e^{-\alpha u}\,du
    \notag\\
    &=
    \left[
    -\frac{1}{\alpha}e^{-\alpha u}
    \right]_0^\infty
    \notag\\
    &=
    \frac{1}{\alpha}.
\end{align}

Hence,

\[
    F(\alpha)=\alpha^{-1}.
\]

By repeated differentiation,

\begin{equation}
    F^{(m)}(\alpha)
    =
    (-1)^m m!\,\alpha^{-(m+1)}.
    \label{eq:gaussian_odd_factorial}
\end{equation}

Substituting into
\eqref{eq:gaussian_odd_recovery},
we obtain

\begin{align}
    I_{2m+1}
    &=
    \frac{(-1)^m}{2}
    \left[
    (-1)^m m!\,\alpha^{-(m+1)}
    \right]_{\alpha=1}
    \notag\\
    &=
    \frac{m!}{2}.
\end{align}

We have therefore established

\begin{equation}
    \boxed{
    \int_0^\infty
    x^{2m+1}e^{-x^2}\,dx
    =
    \frac{m!}{2},
    \qquad m=0,1,2,\ldots
    }
    \label{eq:gaussian_odd_result}
\end{equation}

\medskip

Notice the essential mechanism
behind the calculation. Rather
than attempting to integrate
$u^m e^{-u}$ directly, we
introduced a parameter into
the exponential function.
Differentiation with respect
to that parameter generated
the required power of $u$,
while the parametric integral
itself remained elementary.

This is the principal idea
underlying Feynman's technique.

% -------------------------------------------------
\subsection{Gaussian Integrals: Even Powers}
% -------------------------------------------------

We now turn our attention to Gaussian
integrals involving even powers of $x$.
Our objective is to evaluate

\begin{equation}
    I_{2m}
    =
    \int_0^\infty x^{2m}e^{-x^2}\,dx,
    \qquad m=0,1,2,\ldots
    \label{eq:gaussian_even}
\end{equation}

Unlike the odd-power case, we will
not make an initial substitution.
Instead, we introduce a parameter
directly into the exponent.

\subsubsection*{Step 1: Introducing a parameter}

Consider the parametric integral

\begin{equation}
    G(\alpha)
    =
    \int_0^\infty e^{-\alpha x^2}\,dx,
    \qquad \alpha>0.
    \label{eq:gaussian_even_parameter}
\end{equation}

The reason for introducing $\alpha$
into the exponent is that

\[
    \frac{\partial}{\partial\alpha}
    e^{-\alpha x^2}
    =
    -x^2e^{-\alpha x^2}.
\]

Differentiating again gives

\[
    \frac{\partial^2}{\partial\alpha^2}
    e^{-\alpha x^2}
    =
    x^4e^{-\alpha x^2}.
\]

Continuing this process, we obtain

\begin{equation}
    \frac{\partial^m}{\partial\alpha^m}
    e^{-\alpha x^2}
    =
    (-1)^m x^{2m}e^{-\alpha x^2}.
    \label{eq:gaussian_even_derivative}
\end{equation}

Thus, each differentiation with
respect to $\alpha$ produces an
additional factor of $-x^2$.

Applying Feynman's technique,

\begin{align}
    G^{(m)}(\alpha)
    &=
    \frac{d^m}{d\alpha^m}
    \int_0^\infty e^{-\alpha x^2}\,dx
    \notag\\
    &=
    \int_0^\infty
    \frac{\partial^m}{\partial\alpha^m}
    e^{-\alpha x^2}\,dx
    \notag\\
    &=
    (-1)^m
    \int_0^\infty
    x^{2m}e^{-\alpha x^2}\,dx.
\end{align}

Therefore,

\begin{equation}
    I_{2m}
    =
    (-1)^m
    \left.
    G^{(m)}(\alpha)
    \right|_{\alpha=1}.
    \label{eq:gaussian_even_recovery}
\end{equation}

We have now reduced the problem
to finding the derivatives of
$G(\alpha)$.

\subsubsection*{Step 2: Evaluating the parametric integral}

Recall the classical Gaussian
integral

\begin{equation}
    \int_{-\infty}^{\infty}
    e^{-x^2}\,dx
    =
    \sqrt{\pi}.
    \label{eq:standard_gaussian}
\end{equation}

Since $e^{-x^2}$ is an even
function, it follows that

\begin{equation}
    \int_0^\infty e^{-x^2}\,dx
    =
    \frac{\sqrt{\pi}}{2}.
    \label{eq:half_gaussian}
\end{equation}

We can now evaluate
$G(\alpha)$ by making the
substitution

\[
    u=\sqrt{\alpha}\,x,
    \qquad
    dx=\frac{du}{\sqrt{\alpha}}.
\]

Consequently,

\begin{align}
    G(\alpha)
    &=
    \int_0^\infty
    e^{-\alpha x^2}\,dx
    \notag\\
    &=
    \frac{1}{\sqrt{\alpha}}
    \int_0^\infty e^{-u^2}\,du
    \notag\\
    &=
    \frac{\sqrt{\pi}}{2\sqrt{\alpha}}.
\end{align}

Hence,

\begin{equation}
    \boxed{
    G(\alpha)
    =
    \frac{\sqrt{\pi}}{2}
    \alpha^{-1/2}
    }
    \label{eq:gaussian_even_G}
\end{equation}

\subsubsection*{Step 3: Repeated differentiation}

We now need to compute the
$m$-th derivative of
$\alpha^{-1/2}$.

The first derivative is

\[
    \frac{d}{d\alpha}
    \alpha^{-1/2}
    =
    -\frac{1}{2}\alpha^{-3/2}.
\]

Differentiating a second time,

\[
    \frac{d^2}{d\alpha^2}
    \alpha^{-1/2}
    =
    \frac{1\cdot3}{2^2}
    \alpha^{-5/2}.
\]

The third derivative gives

\[
    \frac{d^3}{d\alpha^3}
    \alpha^{-1/2}
    =
    -\frac{1\cdot3\cdot5}{2^3}
    \alpha^{-7/2}.
\]

Notice the pattern. Every
differentiation introduces
another negative sign and
another odd integer in the
numerator.

After $m$ differentiations,


\begin{equation}
    \frac{d^m}{d\alpha^m}\alpha^{-1/2}
    =
    (-1)^m
    \frac{1\cdot3\cdot5\cdots(2m-1)}{2^m}
    \alpha^{-\left(m+\frac12\right)}.
    \label{eq:half_power_derivative}
\end{equation}


The product

\[
    1\cdot3\cdot5\cdots(2m-1)
\]

is called the \emph{double factorial}
of $2m-1$, denoted by

\[
    (2m-1)!!.
\]

It is important to distinguish
the double factorial from the
ordinary factorial.

For example,

\[
    5!=5\cdot4\cdot3\cdot2\cdot1=120,
\]

whereas

\[
    5!!=5\cdot3\cdot1=15.
\]

Thus, equation
\eqref{eq:half_power_derivative}
can be written as

\begin{equation}
    \frac{d^m}{d\alpha^m}
    \alpha^{-1/2}
    =
    (-1)^m
    \frac{(2m-1)!!}{2^m}
    \alpha^{-(m+1/2)}.
    \label{eq:double_factorial_derivative}
\end{equation}

Substituting into
\eqref{eq:gaussian_even_G},

\begin{equation}
    G^{(m)}(\alpha)
    =
    \frac{\sqrt{\pi}}{2}
    (-1)^m
    \frac{(2m-1)!!}{2^m}
    \alpha^{-(m+1/2)}.
    \label{eq:gaussian_G_derivative}
\end{equation}

Finally, using
\eqref{eq:gaussian_even_recovery}
and setting $\alpha=1$,

\begin{align}
    I_{2m}
    &=
    (-1)^m G^{(m)}(1)
    \notag\\
    &=
    (-1)^m
    \frac{\sqrt{\pi}}{2}
    (-1)^m
    \frac{(2m-1)!!}{2^m}
    \notag\\
    &=
    \frac{\sqrt{\pi}(2m-1)!!}
         {2^{m+1}}.
\end{align}

We have therefore obtained

\begin{equation}
    \boxed{
    \int_0^\infty
    x^{2m}e^{-x^2}\,dx
    =
    \frac{\sqrt{\pi}(2m-1)!!}
         {2^{m+1}}
    }
    \label{eq:gaussian_even_result}
\end{equation}

Here we use the convention
$(-1)!!=1$ when $m=0$,
so that the formula includes
the standard Gaussian integral.

\subsubsection*{Step 4: Expressing the double factorial
in terms of ordinary factorials}

Although the double factorial
arises naturally through repeated
differentiation, it can also be
expressed using ordinary factorials.

To see this, consider

\[
    (2m)!
    =
    1\cdot2\cdot3\cdots(2m).
\]

We can separate this product
into its odd and even factors:

\begin{align}
    (2m)!
    &=
    \underbrace{
    (1\cdot3\cdot5\cdots(2m-1))
    }_{\text{Odd factors}}
    \notag\\
    &\qquad\times
    \underbrace{
    (2\cdot4\cdot6\cdots2m)
    }_{\text{Even factors}}.
\end{align}

The first product is
$(2m-1)!!$.

For the second product,
we can factor out $2$
from each of its $m$ terms:

\begin{align}
    2\cdot4\cdot6\cdots(2m)
    &=
    (2\cdot1)(2\cdot2)
    \cdots(2\cdot m)
    \notag\\
    &=
    2^m(1\cdot2\cdots m)
    \notag\\
    &=
    2^m m!.
\end{align}

Therefore,

\[
    (2m)!
    =
    (2m-1)!!\,2^m m!.
\]

Rearranging,

\begin{equation}
    \boxed{
    (2m-1)!!
    =
    \frac{(2m)!}{2^m m!}
    }
    \label{eq:double_factorial_identity}
\end{equation}

Substituting this identity into
\eqref{eq:gaussian_even_result},
we obtain

\begin{align}
    I_{2m}
    &=
    \frac{\sqrt{\pi}}{2^{m+1}}
    \frac{(2m)!}{2^m m!}
    \notag\\
    &=
    \frac{\sqrt{\pi}(2m)!}
         {2^{2m+1}m!}.
\end{align}

Thus, an equivalent expression
for our Gaussian integral is

\begin{equation}
    \boxed{
    \int_0^\infty
    x^{2m}e^{-x^2}\,dx
    =
    \frac{\sqrt{\pi}(2m)!}
         {2^{2m+1}m!},
    \qquad m=0,1,2,\ldots
    }
    \label{eq:gaussian_even_factorial_result}
\end{equation}

\medskip

We have now evaluated the Gaussian
integrals for both odd and even powers.
For convenience, we summarize the
results:

\begin{equation}
\boxed{
\begin{aligned}
    \int_0^\infty
    x^{2m+1}e^{-x^2}\,dx
    &=
    \frac{m!}{2},\\[6pt]
    \int_0^\infty
    x^{2m}e^{-x^2}\,dx
    &=
    \frac{\sqrt{\pi}(2m)!}
         {2^{2m+1}m!}.
\end{aligned}
}
\label{eq:gaussian_summary}
\end{equation}

The two calculations illustrate
the same underlying principle:
by introducing a parameter into
an exponential function, powers
of the integration variable
can be generated through
repeated differentiation.


% -------------------------------------------------
\subsection{The Dirichlet Integral}
% -------------------------------------------------

We now consider a famous improper integral,
commonly known as the \emph{Dirichlet integral}:

\begin{equation}
    I =
    \int_0^\infty \frac{\sin x}{x}\,dx.
    \label{eq:dirichlet_integral}
\end{equation}

Unlike the Gaussian integrals considered
previously, this integral involves an
oscillatory function divided by $x$.
Moreover, the function $\sin x/x$ does
not have an elementary antiderivative.

Nevertheless, Feynman's technique allows
us to evaluate the definite integral
through a relatively straightforward
calculation.

The central idea is to introduce a
parameter in such a way that
differentiation removes the factor
$1/x$ from the integrand.

\subsubsection*{Step 1: Introducing a parameter}

Consider the parametric integral

\begin{equation}
    F(\alpha)
    =
    \int_0^\infty
    \frac{\sin x}{x}e^{-\alpha x}\,dx,
    \qquad \alpha>0.
    \label{eq:dirichlet_parameter}
\end{equation}

Notice that we have multiplied the
original integrand by an exponential
factor $e^{-\alpha x}$.

This factor serves two purposes.

First, for $\alpha>0$, the exponential
decays rapidly as $x\to\infty$,
making the integral absolutely
convergent.

Second, differentiating
$e^{-\alpha x}$ with respect to
$\alpha$ produces a factor of $-x$,
which will cancel the denominator.

The original integral will eventually
be recovered by allowing
$\alpha\to0^+$.

\subsubsection*{Step 2: Differentiating under the integral sign}

Differentiating $F(\alpha)$,
we obtain

\begin{align}
    F'(\alpha)
    &=
    \frac{d}{d\alpha}
    \int_0^\infty
    \frac{\sin x}{x}e^{-\alpha x}\,dx
    \notag\\
    &=
    \int_0^\infty
    \frac{\partial}{\partial\alpha}
    \left(
    \frac{\sin x}{x}e^{-\alpha x}
    \right)\,dx
    \notag\\
    &=
    \int_0^\infty
    \frac{\sin x}{x}
    (-x)e^{-\alpha x}\,dx.
\end{align}

The factor $x$ in the numerator
cancels with the denominator,
leaving

\begin{equation}
    \boxed{
    F'(\alpha)
    =
    -\int_0^\infty
    e^{-\alpha x}\sin x\,dx.
    }
    \label{eq:dirichlet_derivative}
\end{equation}

This is the crucial simplification.
The original integral contained
the troublesome factor $1/x$,
but differentiation with respect
to the parameter has eliminated it.

We are now left with an elementary
exponential-trigonometric integral.

\subsubsection*{Step 3: Evaluating the resulting integral}

To evaluate the integral on the
right-hand side of
\eqref{eq:dirichlet_derivative},
we use Euler's identity

\[
    e^{ix}=\cos x+i\sin x,
\]

where $i^2=-1$.

In particular,

\[
    \sin x=\operatorname{Im}(e^{ix}),
\]

where $\operatorname{Im}$ denotes
the imaginary part of a complex
number.

Therefore,

\begin{align}
    \int_0^\infty
    e^{-\alpha x}\sin x\,dx
    &=
    \operatorname{Im}
    \left(
    \int_0^\infty
    e^{-\alpha x}e^{ix}\,dx
    \right)
    \notag\\
    &=
    \operatorname{Im}
    \left(
    \int_0^\infty
    e^{(i-\alpha)x}\,dx
    \right).
\end{align}

The integral inside the
imaginary-part operator is
elementary:

\begin{align}
    \int_0^\infty
    e^{(i-\alpha)x}\,dx
    &=
    \left[
    \frac{e^{(i-\alpha)x}}{i-\alpha}
    \right]_0^\infty
    \notag\\
    &=
    -\frac{1}{i-\alpha}
    \notag\\
    &=
    \frac{1}{\alpha-i}.
\end{align}

At this point, it is worth
explaining why the exponential
term vanishes at infinity.

Observe that

\[
    e^{(i-\alpha)x}
    =
    e^{ix}e^{-\alpha x}.
\]

Since

\[
    |e^{ix}|=1,
\]

we have

\begin{equation}
    |e^{(i-\alpha)x}|
    =
    e^{-\alpha x}.
\end{equation}

For $\alpha>0$,

\[
    \lim_{x\to\infty}
    e^{-\alpha x}=0.
\]

Consequently,

\begin{equation}
    \lim_{x\to\infty}
    e^{(i-\alpha)x}=0.
\end{equation}

Thus, although $e^{ix}$ oscillates,
the exponential damping factor
forces the entire expression
to approach zero.

Returning to our calculation,
we obtain

\begin{align}
    \frac{1}{\alpha-i}
    &=
    \frac{\alpha+i}
         {(\alpha-i)(\alpha+i)}
    \notag\\
    &=
    \frac{\alpha+i}{\alpha^2+1}.
\end{align}

Taking the imaginary part gives

\begin{equation}
    \int_0^\infty
    e^{-\alpha x}\sin x\,dx
    =
    \frac{1}{1+\alpha^2}.
\end{equation}

Therefore, equation
\eqref{eq:dirichlet_derivative}
becomes

\begin{equation}
    \boxed{
    F'(\alpha)
    =
    -\frac{1}{1+\alpha^2}.
    }
    \label{eq:dirichlet_differential}
\end{equation}

\subsubsection*{Step 4: Recovering the parametric integral}

We now have a simple differential
equation for $F(\alpha)$:

\[
    \frac{dF}{d\alpha}
    =
    -\frac{1}{1+\alpha^2}.
\]

Integrating with respect to
$\alpha$, we obtain

\begin{equation}
    F(\alpha)
    =
    -\arctan(\alpha)+C,
    \label{eq:dirichlet_antiderivative}
\end{equation}

where $C$ is a constant
of integration.

To determine $C$, we examine
the behavior of $F(\alpha)$
as $\alpha\to\infty$.

Recall that

\[
    F(\alpha)
    =
    \int_0^\infty
    \frac{\sin x}{x}e^{-\alpha x}\,dx.
\]

Since

\[
    \left|\frac{\sin x}{x}\right|
    \leq 1,
    \qquad x>0,
\]

we obtain

\begin{align}
    |F(\alpha)|
    &\leq
    \int_0^\infty
    \left|\frac{\sin x}{x}\right|
    e^{-\alpha x}\,dx
    \notag\\
    &\leq
    \int_0^\infty e^{-\alpha x}\,dx
    \notag\\
    &=
    \frac{1}{\alpha}.
\end{align}

Therefore,

\[
    \lim_{\alpha\to\infty}
    F(\alpha)=0.
\]

Applying this limit to
\eqref{eq:dirichlet_antiderivative},

\[
    0
    =
    -\frac{\pi}{2}+C.
\]

Hence,

\[
    C=\frac{\pi}{2}.
\]

We have consequently obtained

\begin{equation}
    \boxed{
    F(\alpha)
    =
    \frac{\pi}{2}-\arctan(\alpha),
    \qquad \alpha>0.
    }
    \label{eq:dirichlet_parametric_result}
\end{equation}

\subsubsection*{Step 5: Recovering the original integral}

We now return to our original
objective:

\[
    I=
    \int_0^\infty\frac{\sin x}{x}\,dx.
\]

Recall that the parameter
$\alpha$ was introduced only
as an auxiliary variable.

The original integrand is
recovered when the damping
factor $e^{-\alpha x}$
approaches $1$.

Thus, we consider

\[
    \lim_{\alpha\to0^+}F(\alpha).
\]

Using
\eqref{eq:dirichlet_parametric_result},

\begin{align}
    \lim_{\alpha\to0^+}F(\alpha)
    &=
    \lim_{\alpha\to0^+}
    \left(
    \frac{\pi}{2}-\arctan\alpha
    \right)
    \notag\\
    &=
    \frac{\pi}{2}.
\end{align}

The original improper integral
converges by Dirichlet's test.
Moreover, the exponential
damping preserves its value
in the limit $\alpha\to0^+$.
This limiting argument will
be discussed when we examine
the mathematical justification
of the method.

We therefore conclude that

\begin{equation}
    \boxed{
    \int_0^\infty
    \frac{\sin x}{x}\,dx
    =
    \frac{\pi}{2}.
    }
    \label{eq:dirichlet_final}
\end{equation}

\medskip

Notice how the choice of parameter
determined the success of the method.
By introducing an exponential
damping factor, we transformed
an integral involving $\sin x/x$
into an elementary integral
involving $e^{-\alpha x}\sin x$.

In contrast with the Gaussian
examples, where repeated
differentiation generated powers
of the integration variable,
here a single differentiation
eliminated an unwanted factor
from the denominator.

This is another important
application of Feynman's technique.


% -------------------------------------------------
\subsection{A Logarithmic Integral}
% -------------------------------------------------

Our final example concerns a
family of logarithmic integrals.

Consider

\begin{equation}
    I(\alpha)
    =
    \int_0^1
    \frac{x^\alpha-1}{\ln x}\,dx,
    \qquad \alpha>-1.
    \label{eq:logarithmic_integral}
\end{equation}

The presence of $\ln x$
in the denominator makes
direct integration difficult.

Furthermore, the integrand
appears to be undefined at
$x=1$, since both the numerator
and denominator vanish there.

However, this is only a
removable singularity. Indeed,

\[
    \lim_{x\to1}
    \frac{x^\alpha-1}{\ln x}
    =
    \alpha.
\]

Thus, the behavior near
$x=1$ causes no difficulty.

Unlike our previous examples,
the parameter is already
present in the integral.
We therefore do not need to
introduce an additional one.

\subsubsection*{Step 1: Differentiating with respect to the parameter}

We begin by differentiating
$I(\alpha)$:

\begin{align}
    I'(\alpha)
    &=
    \frac{d}{d\alpha}
    \int_0^1
    \frac{x^\alpha-1}{\ln x}\,dx
    \notag\\
    &=
    \int_0^1
    \frac{\partial}{\partial\alpha}
    \left(
    \frac{x^\alpha-1}{\ln x}
    \right)\,dx.
\end{align}

To compute the derivative
inside the integral,
recall that

\[
    x^\alpha=e^{\alpha\ln x}.
\]

Consequently,

\begin{equation}
    \frac{\partial}{\partial\alpha}
    x^\alpha
    =
    x^\alpha\ln x.
\end{equation}

Therefore,

\begin{align}
    I'(\alpha)
    &=
    \int_0^1
    \frac{x^\alpha\ln x}{\ln x}\,dx
    \notag\\
    &=
    \int_0^1 x^\alpha\,dx.
\end{align}

We have obtained

\begin{equation}
    \boxed{
    I'(\alpha)
    =
    \int_0^1 x^\alpha\,dx.
    }
    \label{eq:logarithmic_derivative}
\end{equation}

Notice what happened:
differentiating $x^\alpha$
produced precisely the factor
$\ln x$ required to cancel
the denominator.

This is the essential idea
behind our choice of
differentiation.

\subsubsection*{Step 2: Evaluating the resulting integral}

The integral on the right-hand
side of
\eqref{eq:logarithmic_derivative}
is now elementary.

For $\alpha>-1$,

\begin{align}
    I'(\alpha)
    &=
    \int_0^1 x^\alpha\,dx
    \notag\\
    &=
    \left[
    \frac{x^{\alpha+1}}{\alpha+1}
    \right]_0^1
    \notag\\
    &=
    \frac{1}{\alpha+1}.
\end{align}

Thus,

\begin{equation}
    \boxed{
    I'(\alpha)=\frac{1}{\alpha+1}.
    }
    \label{eq:logarithmic_differential}
\end{equation}

\subsubsection*{Step 3: Recovering the original integral}

We now integrate with respect
to $\alpha$:

\[
    \frac{dI}{d\alpha}
    =
    \frac{1}{\alpha+1}.
\]

Therefore,

\begin{equation}
    I(\alpha)
    =
    \ln(\alpha+1)+C,
    \label{eq:logarithmic_antiderivative}
\end{equation}

where $C$ is a constant
of integration.

To determine this constant,
we use a convenient value
of the parameter.

Setting $\alpha=0$ in
the original integral gives

\begin{align}
    I(0)
    &=
    \int_0^1
    \frac{x^0-1}{\ln x}\,dx
    \notag\\
    &=
    \int_0^1 0\,dx
    \notag\\
    &=0.
\end{align}

On the other hand,
equation
\eqref{eq:logarithmic_antiderivative}
gives

\[
    I(0)=\ln 1+C=C.
\]

Consequently,

\[
    C=0.
\]

We have therefore established

\begin{equation}
    \boxed{
    \int_0^1
    \frac{x^\alpha-1}{\ln x}\,dx
    =
    \ln(\alpha+1),
    \qquad \alpha>-1.
    }
    \label{eq:logarithmic_final}
\end{equation}

As a particular example,
setting $\alpha=1$ gives

\begin{equation}
    \boxed{
    \int_0^1
    \frac{x-1}{\ln x}\,dx
    =
    \ln 2.
    }
    \label{eq:logarithmic_special}
\end{equation}

\medskip

This example is particularly
instructive because the
parameter was already present
in the original integral.

Differentiation generated
a logarithmic factor that
canceled the denominator,
reducing the problem to
an elementary power integral.

The constant of integration
was then determined by
evaluating the original
parametric integral at
a convenient parameter value.


% -------------------------------------------------
\subsection{Observations on the Method}
% -------------------------------------------------

The preceding examples illustrate
several ways in which Feynman's
technique can simplify difficult
integrals.

In the Gaussian integrals,
differentiation with respect to
a parameter generated powers
of the integration variable.

In the Dirichlet integral,
differentiation removed an
unwanted factor of $1/x$.

In the logarithmic integral,
differentiation generated
a factor of $\ln x$ that
canceled the denominator.

Although the calculations
look different, the common
principle remains the same:
we replace a difficult
integration problem with
the study of a suitable
parameter-dependent integral.

Once the derivative with
respect to the parameter
has been evaluated, the
original integral can often
be recovered through
elementary integration
and an appropriate
boundary or limiting condition.

At this point, however,
an important mathematical
question remains unanswered.

Throughout these examples,
we have interchanged
differentiation and integration
without establishing the
conditions under which this
operation is legitimate.

We now turn to the mathematical
foundation of Feynman's technique,
beginning with Leibniz's
integral rule and eventually
examining the problem from
the perspective of linear
operators.


% =================================================
\section{Why Does Feynman's Technique Work?}
% =================================================

In the preceding section, we evaluated
several nontrivial integrals using
Feynman's technique of differentiating
under the integral sign.

Although the calculations were
relatively straightforward, we
repeatedly used an important identity
without establishing its mathematical
validity:

\begin{equation}
    \frac{d}{d\alpha}
    \int_a^b f(x,\alpha)\,dx
    =
    \int_a^b
    \frac{\partial f(x,\alpha)}
         {\partial\alpha}\,dx.
    \label{eq:leibniz_question}
\end{equation}

The left-hand side tells us to
integrate first and then
differentiate, while the
right-hand side reverses
the order of these operations.

Why should these two procedures
produce the same result?

To understand the underlying
idea, we begin with the
definition of integration
through Riemann sums.

% -------------------------------------------------
\subsection{An Intuitive Explanation Using Riemann Sums}
% -------------------------------------------------

Recall from elementary calculus
that a definite integral can be
defined as the limit of a
sequence of Riemann sums.

Suppose that $f(x,\alpha)$
is a sufficiently regular
function on the interval
$[a,b]$.

We divide $[a,b]$ into
$n$ subintervals,

\[
    a=x_0<x_1<\cdots<x_n=b,
\]

and let

\[
    \Delta x_i=x_i-x_{i-1}.
\]

Choosing a point $\xi_i$
from each subinterval,
we construct the Riemann sum

\begin{equation}
    S_n(\alpha)
    =
    \sum_{i=1}^{n}
    f(\xi_i,\alpha)\Delta x_i.
    \label{eq:riemann_parametric}
\end{equation}

As the maximum width of the
subintervals approaches zero,
the Riemann sums converge
to the definite integral:

\begin{equation}
    F(\alpha)
    =
    \int_a^b f(x,\alpha)\,dx
    =
    \lim_{n\to\infty}
    S_n(\alpha).
    \label{eq:riemann_limit}
\end{equation}

Now observe that each
$S_n(\alpha)$ is simply
a finite sum of functions
of the parameter $\alpha$.

Since differentiation is
a linear operation,
we can differentiate
this finite sum term
by term:

\begin{align}
    S_n'(\alpha)
    &=
    \frac{d}{d\alpha}
    \sum_{i=1}^{n}
    f(\xi_i,\alpha)\Delta x_i
    \notag\\
    &=
    \sum_{i=1}^{n}
    \frac{\partial f(\xi_i,\alpha)}
         {\partial\alpha}
    \Delta x_i.
    \label{eq:riemann_derivative}
\end{align}

Notice that $\Delta x_i$
does not depend on $\alpha$,
since the interval of
integration is fixed.

The expression on the
right-hand side of
\eqref{eq:riemann_derivative}
is itself a Riemann sum,
this time for the function

\[
    \frac{\partial f(x,\alpha)}
         {\partial\alpha}.
\]

Consequently, if we take
the limit as the partition
becomes finer, we expect

\begin{equation}
    \lim_{n\to\infty}
    S_n'(\alpha)
    =
    \int_a^b
    \frac{\partial f(x,\alpha)}
         {\partial\alpha}\,dx.
    \label{eq:riemann_derivative_limit}
\end{equation}

On the other hand, since

\[
    F(\alpha)
    =
    \lim_{n\to\infty}S_n(\alpha),
\]

we would like to write

\[
    F'(\alpha)
    =
    \lim_{n\to\infty}S_n'(\alpha).
\]

Combining these observations
would give

\[
    \boxed{
    \frac{d}{d\alpha}
    \int_a^b f(x,\alpha)\,dx
    =
    \int_a^b
    \frac{\partial f(x,\alpha)}
         {\partial\alpha}\,dx.
    }
\]

This is precisely the
identity used in Feynman's
technique.

\medskip

However, there is a subtle
mathematical issue hidden
in this argument.

Although differentiation
can always be distributed
over a finite sum, it
cannot necessarily be
interchanged with an
infinite limit.

In particular, the identity

\[
    \frac{d}{d\alpha}
    \left(
    \lim_{n\to\infty}S_n(\alpha)
    \right)
    =
    \lim_{n\to\infty}S_n'(\alpha)
\]

requires additional
assumptions.

Thus, the Riemann-sum
argument provides a useful
intuition, but it does not,
by itself, constitute
a complete proof.

To establish the result
rigorously, we turn to
Leibniz's integral rule.

% -------------------------------------------------

% -------------------------------------------------
\subsection{Leibniz's Integral Rule}
% -------------------------------------------------

We now establish the mathematical
justification for differentiation
under the integral sign.

We first consider the case where
the limits of integration are fixed,
that is, they do not depend on
the parameter $\alpha$.

\begin{theorem}[Differentiation Under the Integral Sign]

Let $J\subseteq\mathbb{R}$ be an
open interval, and let $a<b$.
Suppose that $f(x,\alpha)$
and its partial derivative

\[
    \frac{\partial f(x,\alpha)}
         {\partial\alpha}
\]

are continuous on $[a,b]\times J$.

Define

\[
    F(\alpha)
    =
    \int_a^b f(x,\alpha)\,dx.
\]

Then $F$ is differentiable
on $J$, and

\begin{equation}
    \boxed{
    F'(\alpha)
    =
    \int_a^b
    \frac{\partial f(x,\alpha)}
         {\partial\alpha}\,dx.
    }
    \label{eq:leibniz_theorem}
\end{equation}

\end{theorem}

\begin{proof}

We will establish the theorem
directly from the definition
of the derivative.

\medskip

\textbf{Step 1: Applying the definition
of the derivative.}

Fix an arbitrary $\alpha\in J$.
By the definition of the derivative,

\begin{equation}
    F'(\alpha)
    =
    \lim_{h\to0}
    \frac{F(\alpha+h)-F(\alpha)}{h}.
\end{equation}

Recall that

\[
    F(\alpha)
    =
    \int_a^b f(x,\alpha)\,dx.
\]

Therefore, for $h\neq0$,

\begin{align}
    &\frac{F(\alpha+h)-F(\alpha)}{h}
    \notag\\
    &=
    \frac{1}{h}
    \left[
    \int_a^b f(x,\alpha+h)\,dx
    -
    \int_a^b f(x,\alpha)\,dx
    \right]
    \notag\\
    &=
    \int_a^b
    \frac{f(x,\alpha+h)-f(x,\alpha)}
         {h}\,dx.
    \label{eq:leibniz_difference_quotient}
\end{align}

The last equality follows
from the linearity of
definite integration.

Notice that, at this stage,
we have not interchanged
differentiation and integration.
We have only used the linearity
of the integral.

Our objective is to prove that

\begin{equation}
    \lim_{h\to0}
    \int_a^b
    \frac{f(x,\alpha+h)-f(x,\alpha)}
         {h}\,dx
    =
    \int_a^b
    \frac{\partial f(x,\alpha)}
         {\partial\alpha}\,dx.
\end{equation}

\medskip

\textbf{Step 2: Applying the Mean Value Theorem.}

For the moment, hold $x$ fixed
and consider $f(x,\alpha)$
as a function of $\alpha$ alone.

Since the partial derivative
with respect to $\alpha$
is continuous, we can apply
the ordinary Mean Value Theorem.

For sufficiently small
$h\neq0$, there exists
a number $\xi_{x,h}$
between $\alpha$ and
$\alpha+h$ such that

\begin{equation}
    \frac{f(x,\alpha+h)-f(x,\alpha)}
         {h}
    =
    \frac{\partial f(x,\xi_{x,h})}
         {\partial\alpha}.
    \label{eq:leibniz_mvt}
\end{equation}

The subscript $x,h$ reminds us
that this intermediate point
may depend on both $x$ and $h$.

Since $\xi_{x,h}$ lies between
$\alpha$ and $\alpha+h$, we have

\begin{equation}
    |\xi_{x,h}-\alpha|
    \leq |h|.
    \label{eq:leibniz_xi_bound}
\end{equation}

Thus, as $h$ approaches zero,
the intermediate point
$\xi_{x,h}$ approaches $\alpha$.

By continuity of the partial
derivative, we expect

\[
    \frac{\partial f(x,\xi_{x,h})}
         {\partial\alpha}
    \longrightarrow
    \frac{\partial f(x,\alpha)}
         {\partial\alpha}.
\]

We must now show that this
convergence is sufficiently
well behaved to pass the limit
through the integral.

\medskip

\textbf{Step 3: Estimating the error.}

Let us compare the difference
quotient of $F$ with the
expression that we want
to obtain.

Define the error

\begin{align}
    E(h)
    &=
    \left|
    \frac{F(\alpha+h)-F(\alpha)}{h}
    -
    \int_a^b
    \frac{\partial f(x,\alpha)}
         {\partial\alpha}\,dx
    \right|.
    \label{eq:leibniz_error}
\end{align}

Using equation
\eqref{eq:leibniz_difference_quotient}
and the linearity of
integration, we obtain

\begin{align}
    E(h)
    &=
    \left|
    \int_a^b
    \left[
    \frac{f(x,\alpha+h)-f(x,\alpha)}{h}
    -
    \frac{\partial f(x,\alpha)}
         {\partial\alpha}
    \right]dx
    \right|.
\end{align}

Now apply the Mean Value
Theorem from Step 2:

\begin{align}
    E(h)
    &=
    \left|
    \int_a^b
    \left[
    \frac{\partial f(x,\xi_{x,h})}
         {\partial\alpha}
    -
    \frac{\partial f(x,\alpha)}
         {\partial\alpha}
    \right]dx
    \right|.
\end{align}

Recall the elementary
inequality

\begin{equation}
    \left|
    \int_a^b g(x)\,dx
    \right|
    \leq
    \int_a^b |g(x)|\,dx.
    \label{eq:integral_absolute_inequality}
\end{equation}

Applying this inequality gives

\begin{align}
    E(h)
    &\leq
    \int_a^b
    \left|
    \frac{\partial f(x,\xi_{x,h})}
         {\partial\alpha}
    -
    \frac{\partial f(x,\alpha)}
         {\partial\alpha}
    \right|dx.
    \label{eq:leibniz_error_bound}
\end{align}

We now want to show that
the quantity inside the
absolute value becomes
uniformly small as
$h\to0$.

\medskip

\textbf{Step 4: Why the error becomes small.}

Since $J$ is open and
$\alpha\in J$, we can
choose a small number
$\delta>0$ such that

\[
    [\alpha-\delta,\alpha+\delta]
    \subset J.
\]

Consider the rectangle

\[
    K=
    [a,b]\times
    [\alpha-\delta,\alpha+\delta].
\]

This rectangle is closed
and bounded, hence compact.

By assumption, the partial
derivative

\[
    \frac{\partial f(x,t)}
         {\partial\alpha}
\]

is continuous on $K$.

A standard result from
elementary analysis states
that a continuous function
on a compact set is
uniformly continuous.

In our setting, this means
that for every $\varepsilon>0$,
we can choose a sufficiently
small $\eta>0$ such that

\begin{equation}
    |t-\alpha|<\eta
    \quad\Longrightarrow\quad
    \left|
    \frac{\partial f(x,t)}
         {\partial\alpha}
    -
    \frac{\partial f(x,\alpha)}
         {\partial\alpha}
    \right|
    <\varepsilon
    \label{eq:leibniz_uniform_continuity}
\end{equation}

for every $x\in[a,b]$,
provided both parameter
values lie in the rectangle.

The important point is that
we can use the same $\eta$
for every $x\in[a,b]$.

This is precisely what
uniform continuity guarantees.

Now recall from
\eqref{eq:leibniz_xi_bound}
that

\[
    |\xi_{x,h}-\alpha|
    \leq |h|.
\]

Therefore, whenever $h$
is sufficiently small,
equation
\eqref{eq:leibniz_uniform_continuity}
implies

\begin{equation}
    \left|
    \frac{\partial f(x,\xi_{x,h})}
         {\partial\alpha}
    -
    \frac{\partial f(x,\alpha)}
         {\partial\alpha}
    \right|
    <\varepsilon
    \label{eq:leibniz_small_integrand}
\end{equation}

for every $x\in[a,b]$.

Substituting this bound
into equation
\eqref{eq:leibniz_error_bound},
we obtain

\begin{align}
    E(h)
    &\leq
    \int_a^b
    \left|
    \frac{\partial f(x,\xi_{x,h})}
         {\partial\alpha}
    -
    \frac{\partial f(x,\alpha)}
         {\partial\alpha}
    \right|dx
    \notag\\
    &\leq
    \int_a^b \varepsilon\,dx
    \notag\\
    &=
    \varepsilon(b-a).
\end{align}



Geometrically, this is
the area of a rectangle
with height $\varepsilon$
and width $(b-a)$.

Thus, we have established

\begin{equation}
    \boxed{
    0\leq E(h)\leq\varepsilon(b-a)
    }
    \label{eq:leibniz_final_error_bound}
\end{equation}

for all sufficiently
small nonzero $h$.

Since $\varepsilon>0$
can be chosen arbitrarily
small and $(b-a)$ is
a fixed finite number,
the error must approach zero.

More precisely, given any
desired error tolerance
$\rho>0$, we may choose

\[
    \varepsilon=\frac{\rho}{b-a}.
\]

Then the preceding
inequality gives

\[
    0\leq E(h)\leq\rho
\]

whenever $h$ is
sufficiently small.

Consequently,

\begin{equation}
    \lim_{h\to0}E(h)=0.
\end{equation}

\medskip

\textbf{Step 5: Completing the proof.}

By the definition of
$E(h)$, the preceding
limit means that

\[
    \lim_{h\to0}
    \frac{F(\alpha+h)-F(\alpha)}{h}
    =
    \int_a^b
    \frac{\partial f(x,\alpha)}
         {\partial\alpha}\,dx.
\]

But the expression on
the left-hand side is
exactly $F'(\alpha)$.

Therefore,

\begin{equation}
    \boxed{
    \frac{d}{d\alpha}
    \int_a^b f(x,\alpha)\,dx
    =
    \int_a^b
    \frac{\partial f(x,\alpha)}
         {\partial\alpha}\,dx.
    }
\end{equation}

Since $\alpha\in J$ was
arbitrary, the result
holds throughout $J$.

This completes the proof.

\end{proof}

\medskip

The proof reveals the
mathematical principle
behind Feynman's technique.



% -------------------------------------------------
\subsection{When the Limits of Integration Depend on the Parameter}
% -------------------------------------------------

The preceding theorem assumes
that the integration limits
$a$ and $b$ are fixed.

Suppose instead that the
limits are differentiable
functions of $\alpha$:

\[
    F(\alpha)
    =
    \int_{a(\alpha)}^{b(\alpha)}
    f(x,\alpha)\,dx.
\]

Under suitable continuity
and differentiability
assumptions, the general
form of Leibniz's rule is

\begin{equation}
\boxed{
\begin{aligned}
    F'(\alpha)
    &=
    f(b(\alpha),\alpha)b'(\alpha)\\
    &\quad-
    f(a(\alpha),\alpha)a'(\alpha)\\
    &\quad+
    \int_{a(\alpha)}^{b(\alpha)}
    \frac{\partial f(x,\alpha)}
         {\partial\alpha}\,dx.
\end{aligned}
}
\label{eq:general_leibniz_rule}
\end{equation}

The first two terms account
for the movement of the
integration boundaries,
while the third term
accounts for the change
in the integrand itself.

When $a$ and $b$ are
constants,

\[
    a'(\alpha)=b'(\alpha)=0,
\]

and the formula reduces
to the familiar identity

\[
    \frac{d}{d\alpha}
    \int_a^b f(x,\alpha)\,dx
    =
    \int_a^b
    \frac{\partial f(x,\alpha)}
         {\partial\alpha}\,dx.
\]

This is the special case
we have used throughout
our examples.

% -------------------------------------------------
\subsection{A Remark on Improper Integrals}
% -------------------------------------------------

There is one more subtlety
worth mentioning.

Several integrals considered
in the preceding section
were improper integrals,
with an infinite integration
interval or a singularity
at an endpoint.

The theorem proved above
applies directly to a
finite interval $[a,b]$.
For an improper integral
such as

\[
    F(\alpha)
    =
    \int_0^\infty f(x,\alpha)\,dx,
\]

additional conditions are
required to justify
differentiating under the
integral sign.

One useful sufficient
condition is that, for
parameters in a neighborhood
of a fixed value $\alpha_0$,
there exists an integrable
function $g(x)$ such that

\begin{equation}
    \left|
    \frac{\partial f(x,\alpha)}
         {\partial\alpha}
    \right|
    \leq g(x),
    \label{eq:dominating_function}
\end{equation}

where

\[
    \int_0^\infty g(x)\,dx<\infty.
\]

If the integrand is
differentiable with respect
to $\alpha$, the integral
defining $F(\alpha)$ is
well-defined, and the
necessary measurability
conditions hold, this
domination allows us to
justify the interchange
using the dominated
convergence theorem.

For example, the Gaussian
integrands we studied are
smooth and decay rapidly
at infinity whenever
$\alpha>0$.

Their parameter derivatives
can be controlled by
integrable functions on
suitable neighborhoods
of any fixed positive
parameter.

The exponentially damped
Dirichlet integral admits
a similar justification
for $\alpha>0$.

However, recovering the
original Dirichlet integral
by letting $\alpha\to0^+$
is a separate limiting
question. Its justification
cannot be obtained merely
by substituting $\alpha=0$
into the differentiation
formula.

These distinctions are
important because they
show that Feynman's technique
is not an unrestricted
algebraic manipulation.
Its validity depends on
the properties of the
functions and the domains
under consideration.

\medskip

We have now obtained the
mathematical foundation
for differentiating under
the integral sign.

Nevertheless, another
interesting interpretation
remains.

Since both integration
and differentiation are
linear operations, we may
ask whether Feynman's
technique can be understood
in the language of
linear operators.

This perspective will
be developed in the
next section.

% =================================================
\section{An Operator-Theoretic Interpretation
of Feynman's Technique}
% =================================================

We have now established the mathematical
conditions under which differentiation
and integration can be interchanged. However, there is another interesting
way to understand Feynman's technique,
particularly for readers familiar
with linear algebra and matrix mechanics.

Instead of treating differentiation
and integration merely as operations
from calculus, we can regard them as
linear operators acting on spaces
of functions.

From this perspective, the central
question becomes:

\medskip

\emph{Under what conditions do the
differentiation and integration
operators commute?}

\medskip

We will develop this interpretation
step by step.

% -------------------------------------------------
\subsection{Functions as Vectors}
% -------------------------------------------------

Recall that a vector space is a
collection of objects on which
vector addition and scalar
multiplication are defined,
satisfying the usual vector
space axioms.

In elementary linear algebra,
we generally work with vectors
in $\mathbb{R}^n$.

For example,

\[
    v=
    \begin{pmatrix}
        v_1\\
        v_2\\
        \vdots\\
        v_n
    \end{pmatrix}
    \in\mathbb{R}^n.
\]

However, vectors need not be
finite lists of real numbers.

Functions can also be regarded
as vectors.

For example, the collection
$C([a,b])$ of continuous
real-valued functions on
$[a,b]$ forms a vector space
under the operations

\[
    (f+g)(x)=f(x)+g(x)
\]

and

\[
    (cf)(x)=cf(x),
    \qquad c\in\mathbb{R}.
\]

Unlike $\mathbb{R}^n$,
this function space is
infinite-dimensional.

The important observation
is that familiar operations
from calculus can be studied
as transformations between
such vector spaces.

In particular, differentiation
and integration are linear
transformations.

% -------------------------------------------------
\subsection{Differentiation and Integration
as Linear Operators}
% -------------------------------------------------

Let us first consider the
differentiation operator

\[
    Df=\frac{df}{dx}.
\]

For differentiable functions
$f$ and $g$, and constants
$c_1,c_2$, we have

\begin{align}
    D(c_1f+c_2g)
    &=
    \frac{d}{dx}(c_1f+c_2g)
    \notag\\
    &=
    c_1\frac{df}{dx}
    +
    c_2\frac{dg}{dx}
    \notag\\
    &=
    c_1Df+c_2Dg.
\end{align}

Thus, differentiation is
a linear operator.

Similarly, define the
integration operator

\begin{equation}
    Tf=\int_a^b f(x)\,dx.
\end{equation}

By linearity of integration,

\begin{align}
    T(c_1f+c_2g)
    &=
    \int_a^b
    (c_1f(x)+c_2g(x))\,dx
    \notag\\
    &=
    c_1\int_a^b f(x)\,dx
    +
    c_2\int_a^b g(x)\,dx
    \notag\\
    &=
    c_1Tf+c_2Tg.
\end{align}

Therefore, integration is
also a linear operator.

Notice that differentiation
and definite integration
need not have the same
domain and codomain.

For example, differentiation
maps differentiable functions
to functions, whereas
definite integration over
a fixed interval maps
functions to real numbers.

When a parameter is present,
the situation becomes
particularly interesting.

% -------------------------------------------------
\subsection{Introducing the Parameter}
% -------------------------------------------------

Suppose that we have a
function $f(x,\alpha)$,
where $x$ is the variable
of integration and
$\alpha$ is a parameter.

Let $J$ be an open interval
containing the allowed
values of $\alpha$.

We define the integration
operator with respect to
$x$ by

\begin{equation}
    (T_xf)(\alpha)
    =
    \int_a^b f(x,\alpha)\,dx.
    \label{eq:operator_T}
\end{equation}

Notice that $T_xf$ is
now a function of $\alpha$.

Integration has eliminated
the dependence on $x$,
but the parameter remains.

Next, define differentiation
with respect to $\alpha$:

\begin{equation}
    (D_\alpha f)(x,\alpha)
    =
    \frac{\partial f(x,\alpha)}
         {\partial\alpha}.
    \label{eq:operator_D}
\end{equation}

We also use $D_\alpha$
to denote ordinary
differentiation when
acting on a function
of $\alpha$ alone.

Both operations are linear,
but they act in different
ways.

The integration operator
$T_x$ integrates over
the variable $x$,
whereas $D_\alpha$
measures variation with
respect to the parameter
$\alpha$.

Our objective is to
investigate whether
the order of these
operations matters.

% -------------------------------------------------
\subsection{Commutativity of the Operators}
% -------------------------------------------------

Consider first applying
the integration operator
and then differentiating
the result.

We obtain

\begin{equation}
    D_\alpha(T_xf)
    =
    \frac{d}{d\alpha}
    \left(
    \int_a^b f(x,\alpha)\,dx
    \right).
    \label{eq:operator_DT}
\end{equation}

Now reverse the order.

First differentiate the
integrand with respect
to $\alpha$, and then
integrate over $x$:

\begin{equation}
    T_x(D_\alpha f)
    =
    \int_a^b
    \frac{\partial f(x,\alpha)}
         {\partial\alpha}\,dx.
    \label{eq:operator_TD}
\end{equation}

Feynman's technique asks
whether these two
expressions are equal.

In other words,

\begin{equation}
    \boxed{
    D_\alpha(T_xf)
    =
    T_x(D_\alpha f).
    }
    \label{eq:operator_commutation}
\end{equation}

But this is exactly
the statement of
Leibniz's integral rule
for fixed integration limits.

Consequently, under
the conditions established
in the preceding section,
the two operations can
be performed in either
order without changing
the result.

This provides an
operator-theoretic
interpretation of
differentiation under
the integral sign.

\subsubsection*{The commutator}

In linear algebra and
operator theory, the
\emph{commutator} of two
operators $A$ and $B$
acting on a common
vector space is defined by

\begin{equation}
    [A,B]=AB-BA.
    \label{eq:commutator_definition}
\end{equation}

The operators commute
precisely when

\[
    [A,B]=0.
\]

We would therefore like
to express Feynman's
technique in the form

\[
    [D_\alpha,T_x]=0.
\]

There is a small technical
point here.

The integration operator
defined in
\eqref{eq:operator_T}
maps a function of
$(x,\alpha)$ to a function
of $\alpha$ alone.

Strictly speaking,
$D_\alpha$ and $T_x$
are therefore not both
endomorphisms of the
same function space.

The precise mathematical
statement is the identity
\eqref{eq:operator_commutation},
where differentiation acts
on the appropriate space
before or after integration.

If we wish to use the
ordinary commutator notation,
we can regard the output
of $T_x$ as a function
that is constant in $x$.

More explicitly, define

\begin{equation}
    (\widetilde{T}_x f)(x,\alpha)
    =
    \int_a^b f(t,\alpha)\,dt.
    \label{eq:extended_integration_operator}
\end{equation}

Here $t$ is merely a
dummy integration variable.

Unlike the original
integration operator,
$\widetilde{T}_x$ maps
functions of $(x,\alpha)$
back into functions of
$(x,\alpha)$.

We can now form the
commutator in the
usual sense:

\begin{equation}
    [D_\alpha,\widetilde{T}_x]
    =
    D_\alpha\widetilde{T}_x
    -
    \widetilde{T}_xD_\alpha.
    \label{eq:feynman_commutator}
\end{equation}

Under the hypotheses
of Leibniz's rule,

\begin{equation}
    \boxed{
    [D_\alpha,\widetilde{T}_x]f=0.
    }
    \label{eq:vanishing_commutator}
\end{equation}

Thus, on a suitable
space of functions,
the differentiation
and fixed-interval
integration operators
commute.

This is the operator
identity underlying
Feynman's technique.

% -------------------------------------------------
\subsection{A Finite-Dimensional Analogy}
% -------------------------------------------------

We can make the preceding
discussion more intuitive
by considering an analogous
problem in ordinary
linear algebra.

Suppose that

\[
    v(\alpha)=
    \begin{pmatrix}
        v_1(\alpha)\\
        v_2(\alpha)\\
        \vdots\\
        v_n(\alpha)
    \end{pmatrix}
\]

is a differentiable
vector-valued function.

Let

\[
    w=
    \begin{pmatrix}
        w_1\\
        w_2\\
        \vdots\\
        w_n
    \end{pmatrix}
\]

be a fixed vector,
independent of $\alpha$.

Consider the scalar-valued
function

\[
    F(\alpha)=w^Tv(\alpha).
\]

By the definition of
matrix multiplication,

\[
    F(\alpha)
    =
    \sum_{i=1}^{n}
    w_i v_i(\alpha).
\]

Differentiating with
respect to $\alpha$,

\begin{align}
    F'(\alpha)
    &=
    \frac{d}{d\alpha}
    \sum_{i=1}^{n}
    w_i v_i(\alpha)
    \notag\\
    &=
    \sum_{i=1}^{n}
    w_i v_i'(\alpha)
    \notag\\
    &=
    w^Tv'(\alpha).
\end{align}

Therefore,

\begin{equation}
    \boxed{
    \frac{d}{d\alpha}
    \bigl(w^Tv(\alpha)\bigr)
    =
    w^Tv'(\alpha).
    }
    \label{eq:finite_dimensional_analogy}
\end{equation}

Notice that differentiation
can be performed before
or after taking the
weighted sum.

Now recall that a definite
integral can be approximated
by a Riemann sum:

\[
    \int_a^b f(x,\alpha)\,dx
    \approx
    \sum_{i=1}^{n}
    f(\xi_i,\alpha)\Delta x_i.
\]

This is precisely a
weighted sum of the form

\[
    w^Tv(\alpha),
\]

where

\[
    w_i=\Delta x_i,
    \qquad
    v_i(\alpha)=f(\xi_i,\alpha).
\]

Since the weights
$\Delta x_i$ are
independent of $\alpha$,
we can differentiate
the finite sum
componentwise.

The passage from finite
sums to integrals is
what makes Feynman's
technique a natural
infinite-dimensional
analogue of
\eqref{eq:finite_dimensional_analogy}.

However, unlike a finite
sum, an integral involves
a limiting process.

Consequently, additional
regularity conditions
are needed to justify
the interchange in the
infinite-dimensional case.

% -------------------------------------------------
\subsection{Why Linearity Alone Is Not Sufficient}
% -------------------------------------------------

It is important to emphasize
that the vanishing commutator
does not follow merely from
the linearity of the two
operators.

Indeed, even in ordinary
linear algebra, two matrices
$A$ and $B$ may both represent
linear transformations
without commuting:

\[
    AB\neq BA.
\]

Similarly, differentiation
and integration are linear
operations, but this fact
alone does not establish
that their order can
always be interchanged.

The equality

\[
    D_\alpha T_xf
    =
    T_xD_\alpha f
\]

must be justified by
the properties of the
function $f$ and the
domain on which the
operators act.

For continuously
differentiable integrands
over fixed finite
intervals, Leibniz's rule
provides that justification.

For improper integrals,
additional convergence
conditions are required.

Furthermore, if the
integration limits depend
on the parameter, then
Leibniz's general rule gives

\begin{align}
    &\frac{d}{d\alpha}
    \int_{a(\alpha)}^{b(\alpha)}
    f(x,\alpha)\,dx
    \notag\\
    &\quad-
    \int_{a(\alpha)}^{b(\alpha)}
    \frac{\partial f(x,\alpha)}
         {\partial\alpha}\,dx
    \notag\\
    &=
    f(b(\alpha),\alpha)b'(\alpha)
    -
    f(a(\alpha),\alpha)a'(\alpha).
\end{align}

Thus, additional boundary
terms appear when the
limits move with the
parameter.

The operator identity
therefore has a precise
domain of validity, and one should be careful before exploiting it in scientific computing.

% -------------------------------------------------
\subsection{Returning to Feynman's Technique}
% -------------------------------------------------

Let us briefly revisit
the Gaussian integral
from the operator perspective.

Recall that

\begin{equation}
    G(\alpha)
    =
    \int_0^\infty
    e^{-\alpha x^2}\,dx
    =
    \frac{\sqrt{\pi}}
         {2\sqrt{\alpha}},
    \qquad \alpha>0.
\end{equation}

Differentiating the
explicit expression,

\begin{equation}
    G'(\alpha)
    =
    -\frac{\sqrt{\pi}}
          {4\alpha^{3/2}}.
\end{equation}

On the other hand,
differentiating under
the integral sign,

\begin{align}
    G'(\alpha)
    &=
    \int_0^\infty
    \frac{\partial}{\partial\alpha}
    e^{-\alpha x^2}\,dx
    \notag\\
    &=
    -\int_0^\infty
    x^2e^{-\alpha x^2}\,dx.
\end{align}

Equating the two
expressions gives

\[
    \int_0^\infty
    x^2e^{-\alpha x^2}\,dx
    =
    \frac{\sqrt{\pi}}
         {4\alpha^{3/2}}.
\]

At $\alpha=1$, we obtain

\begin{equation}
    \boxed{
    \int_0^\infty
    x^2e^{-x^2}\,dx
    =
    \frac{\sqrt{\pi}}{4}.
    }
\end{equation}

In operator notation,
the essential step was

\[
    D_\alpha(T_xf)
    =
    T_x(D_\alpha f),
\]

where

\[
    f(x,\alpha)=e^{-\alpha x^2}.
\]

The two operators
produced the same result,
but one order of
calculation was much
more useful than
the other.

This is precisely
the computational
advantage exploited
by Feynman's technique.







% =================================================
\section{Conclusion}
% =================================================

In this note, we have explored Feynman's
technique of parametric integration,
beginning with its computational applications
and concluding with its mathematical
justification.

We first considered several examples
involving Gaussian, trigonometric, and
logarithmic integrals. In each case,
the essential idea was to introduce
or exploit a parameter whose
differentiation simplifies the
original integrand.

The Gaussian integrals illustrated
how repeated differentiation
generates powers of the integration
variable and naturally leads to
ordinary and double factorials.
The Dirichlet integral demonstrated
how exponential damping can
transform a difficult oscillatory
integral into an elementary one.
Finally, the logarithmic integral
showed how differentiation can
eliminate an otherwise troublesome
logarithmic denominator.

Having developed the computational
method, we investigated why
differentiation and integration
can sometimes be interchanged.
Starting from the Riemann-sum
interpretation of integration,
we were led to Leibniz's
integral rule, which provides
a rigorous mathematical
foundation for the technique.

We then examined the method
from the perspective of
linear algebra and operator
theory. By regarding
differentiation and integration
as linear operators acting
on suitable spaces of functions,
we interpreted Leibniz's rule
as a commutation relation.

In particular, under appropriate
regularity and convergence
conditions, we have

\begin{equation}
    D_\alpha T_x f
    =
    T_x D_\alpha f.
\end{equation}

Equivalently, using the
extended integration operator
introduced earlier, the
corresponding commutator
vanishes:

\begin{equation}
    \boxed{
    [D_\alpha,\widetilde{T}_x]f=0.
    }
\end{equation}

It is important to remember,
however, that this relation
does not hold for arbitrary
functions. Its validity
depends on the mathematical
conditions governing the
interchange of differentiation
and integration.

The broader lesson is that
many difficult mathematical
problems become simpler when
we recognize the underlying
structure of the operations
involved. Rather than
evaluating an integral
directly, Feynman's technique
allows us to embed it in
a family of integrals
whose dependence on a
parameter is easier
to investigate.

Thus, what initially appears
to be a clever computational
trick is, in fact, an
application of a deeper
principle from mathematical
analysis and operator theory.

The technique illustrates
an important theme in
mathematics: understanding
when and why operations
can be interchanged is
often just as valuable
as knowing how to
perform them.


% =================================================
\begin{thebibliography}{9}
% =================================================



\bibitem{rudin}
Walter Rudin,
\textit{Principles of Mathematical Analysis},
3rd edition,
McGraw-Hill, 1976.


\bibitem{boas}
Mary L. Boas,
\textit{Mathematical Methods in the Physical Sciences},
3rd edition,
John Wiley \& Sons, 2005.


\bibitem{freericks}
James K. Freericks,
\textit{Math Methods: Physics 155},
YouTube lecture series.
Available at:
\url{https://www.youtube.com/playlist?list=PLxyBMSdBhYxMsgxc2t83yyc3g6AM8GY6D}.


\end{thebibliography}

\end{document}

\end{document}
